% Part: computability
% Chapter: tm-computations
% Section: church-turing-thesis

\documentclass[../../../include/open-logic-section]{subfiles}

\begin{document}

% SEGMENT OLP-0263-S01

\olfileid{tur}{mac}{ctt}
\olsection{சர்ச்--டியூரிங் முன்வைப்பு}

விளைவுறு செயல்முறை என்ற கருத்துக்கான ஒரு துல்லியமான மாற்றீடாக
டியூரிங் பொறிகள் அமைய வேண்டும். விளைவுறு செயல்முறை என்ற கருத்தையும்
டியூரிங் பொறி என்ற கருத்தையும் புரிந்துகொண்ட எவருக்கும், ஒரு விளைவுறு
செயல்முறையால் செய்யக்கூடிய எதையும் ஒரு டியூரிங் பொறியால் செய்ய முடியும்
என்ற உள்ளுணர்வு ஏற்படும் என்று டியூரிங் கருதினார். விளைவுறு செயல்முறை
என்ற கருத்துக்காக முன்மொழியப்பட்ட மற்ற எல்லாத் துல்லியமான மாற்றீடுகளும்
டியூரிங் பொறி என்ற கருத்துக்கு விரிவளவில் சமானமானவையாக அமைகின்றன
என்ற உண்மை இக்கூற்றை ஆதரிக்கிறது---அதாவது, அவற்றால் சரியாக அதே
சார்புகளின் கணத்தைக் கணிக்க முடியும். இக்கூற்று
\emph{சர்ச்--டியூரிங் முன்வைப்பு} எனப்படுகிறது.

\begin{defn}[சர்ச்--டியூரிங் முன்வைப்பு]
\emph{சர்ச்--டியூரிங் முன்வைப்பு}, ஒரு விளைவுறு செயல்முறையால்
கணிக்கத்தக்க எதுவும் டியூரிங்-கணிக்கத்தக்கது என்று கூறுகிறது.
\end{defn}

சர்ச்--டியூரிங் முன்வைப்பு இரு வழிகளில் பயன்படுத்தப்படுகிறது. அதன்
முதல் வகைப் பயன்பாடு சோம்பலுக்கான ஒரு சாக்கு. ஏதோ ஒன்றைக் கணிப்பதற்கான
ஒரு விளைவுறு செயல்முறையின் விவரிப்பு, எடுத்துக்காட்டாக
``போலிக்குறிமுறையில்'', நம்மிடம் உள்ளது என்று கொள்க. நாம் உண்மையில்
அதை உருவாக்கியிருக்காவிட்டாலும், ஏதோ ஒரு டியூரிங் பொறி அதே சார்பைக்
கணிக்கிறது என்ற கூற்றை நியாயப்படுத்த சர்ச்--டியூரிங் முன்வைப்பைப்
பயன்படுத்தலாம்.

சர்ச்--டியூரிங் முன்வைப்பின் மற்றொரு பயன்பாடு மெய்யியல் நோக்கில் மேலும்
ஆர்வமூட்டுகிறது. டியூரிங் பொறிகளால் கணிக்க முடியாத சார்புகள் உள்ளன
என்று காட்டலாம். இதிலிருந்து சர்ச்--டியூரிங் முன்வைப்பைப் பயன்படுத்தி,
எந்தச் செயல்முறையைப் பயன்படுத்தினாலும் அதை விளைவுறக் கணிக்க முடியாது
என்று முடிவுசெய்யலாம். ஏனெனில், அத்தகைய ஒரு செயல்முறை இருந்தால்,
சர்ச்--டியூரிங் முன்வைப்பின்படி அதற்கான ஒரு டியூரிங் பொறியும் இருக்கும்.
எனவே அதைக் கணிக்கும் டியூரிங் பொறி இல்லை என்று நிறுவ முடிந்தால்,
விளைவுறு செயல்முறையும் இருக்க முடியாது. குறிப்பாக, நிற்றல் சிக்கல்
எனப்படுவது டியூரிங் பொறிகளால் தீர்க்கப்பட முடியாது என்பதோடு, அதை
எந்த வகையிலும் விளைவுறத் தீர்க்க முடியாது என்றும் கூறுவதற்குச்
சர்ச்--டியூரிங் முன்வைப்பு பயன்படுத்தப்படுகிறது.

\end{document}
